\chapter{Conclusion}
\label{ch:conclusion}

This dissertation has investigated the political economy of formal-informal linkages in India's textile and apparel manufacturing, adjudicating between the modernization-driven displacement narrative and the structuralist adverse incorporation narrative. By integrating formal sector ASI data with informal sector ASUSE data into a state-industry panel, and triangulating empirical findings with theoretical simulations, the study provides a comprehensive account of how labor institutions shape surplus distribution in contemporary value chains.

\section{Summary of Findings}
\label{sec:summary_findings}

The research yields four principal conclusions:

First, formal sector productivity growth is a primary driver of \textbf{strategic outsourcing} to the informal sector. This relationship is not uniform but is highly sensitive to the regulatory environment: states with stricter labor enforcement exhibit a 4.4 times larger outsourcing elasticity. This suggests that informality is strategically deployed as a mechanism for regulatory arbitrage rather than being a transitory byproduct of low productivity.

Second, there is a complete \textbf{absence of productivity pass-through} to informal wages. Across six robustness specifications, the elasticity of informal wages with respect to formal productivity remains statistically zero. This confirms a monopsonistic wage-setting regime where formal lead firms capture the entire surplus created by productivity gains, leaving informal workers at near-subsistence earnings regardless of the value they generate.

Third, the informal sector exhibits \textbf{extreme structural persistence}. With an autoregressive coefficient of 0.90, informal employment levels are almost entirely predetermined by their own history and are largely insensitive to contemporaneous formal productivity shocks. This path dependence characterizes the Indian textile informal sector as a self-perpetuating structural feature of the economy rather than a residual category awaiting transition.

Fourth, \textbf{policy simulations validate the adverse incorporation model}. The simulation shows that conventional labor enforcement merely scales the informal sector (industrial density) without improving its terms of incorporation (wages). The critical divergence—where employment responds to enforcement while wages remain flat—is the hallmark of an adverse incorporation equilibrium.

\section{Contributions of the Study}
\label{sec:contributions}

This research makes three key contributions:

\textbf{1. Theoretical Contribution:} It develops a tractable \textbf{CES-Nash bargaining framework} that nests both displacement and adverse incorporation as competing parameter restrictions. By mapping qualitative industrial narratives to estimable parameters, the model provides a rigorous foundation for future structural research on dual economies.

\textbf{2. Empirical Contribution:} It provides the \textbf{first quantitative evidence} on the "price squeeze" mechanism in India's manufacturing linkages. By identifying the zero pass-through of productivity to wages, the study moves beyond descriptive accounts of the wage gap to isolate the distributional mechanism at the heart of informality.

\textbf{3. Policy Contribution:} It identifies a \textbf{"regulatory arbitrage trap"} in India's labor governance. The finding that enforcement perversely incentivizes outsourcing to the unregulated informal sector suggests that traditional site-based regulation must be replaced by value-chain-wide accountability frameworks.

\section{Directions for Future Research}
\label{sec:future_directions}

Three directions for future research follow naturally from these conclusions. First, the use of \textbf{restricted-access ASI microdata} would allow for district-level analysis, providing sharper identification of local labor market effects and addressing the geographic suppression problem. Second, the integration of \textbf{GST-linked transaction data} could provide a direct firm-to-firm test of the price-squeeze mechanism, sidestepping the aggregation issues of survey data. Third, \textbf{longitudinal tracking of informal enterprises} would allow researchers to track "graduation" versus "lock-in" dynamics over time, providing a definitive temporal resolution to the dualism debate.

In conclusion, the persistence of informality in India's textile sector is not a sign of modernization's failure, but of its strategic success. As formal firms grow more productive, they do not dissolve the informal sector; they deepen its adverse incorporation. Addressing this requires a fundamental shift in labor policy—from regulating the factory gate to governing the global value chain.
